\documentclass[a4paper]{article}

%Tutti gli usepackage vanno qui

\usepackage{geometry}
\usepackage[italian]{babel}
\usepackage[utf8]{inputenc}
\usepackage[T1]{fontenc}
\usepackage{tgschola}
%\usepackage{tgbonum}
\usepackage{tabularx}
\usepackage{longtable}
\usepackage{hyperref}
\usepackage{enumitem}
\usepackage[toc]{appendix}
\hypersetup{
	colorlinks=true,
	linkcolor=black,
	filecolor=magenta,      
	urlcolor=blue,
}
% Numerazione figure 
\let\counterwithout\relax
\let\counterwithin\relax
\usepackage{chngcntr}

\counterwithin{table}{subsection}
\counterwithin{figure}{subsection}

\usepackage[bottom]{footmisc}
\usepackage{fancyhdr}
\setcounter{secnumdepth}{4}
\usepackage{amsmath, amssymb}
\usepackage{array}
\usepackage{graphicx}

\usepackage{ifthen}

%\usepackage{float}
\usepackage{layouts}
\usepackage{url}
\usepackage{comment}
\usepackage{float}
\usepackage{eurosym}

\usepackage{lastpage}
\usepackage{layouts}
\usepackage{float}
\usepackage{eurosym}

%Comandi di impaginazione uguale per tutti i documenti
\pagestyle{fancy}
\lhead{\includegraphics[scale=0.04]{../../../../latex/images/logoTWM.png}}
%Titolo del documento
\rhead{\doctitle{}}
%\rfoot{\thepage}
\cfoot{Pagina \thepage\ di \pageref{LastPage}}
\setlength{\headheight}{35pt}
\setcounter{tocdepth}{5}
\setcounter{secnumdepth}{5}
\renewcommand{\footrulewidth}{0.4pt}

% multirow per tabelle
\usepackage{multirow}

% Permette tabelle su più pagine
%\usepackage{longtable}


% colore di sfondo per le celle
\usepackage[table]{xcolor}

%COMANDI TABELLE
\newcommand{\rowcolorhead}{\rowcolor[HTML]{9b240a}} %intestazione 
% check for missing commands
\newcommand{\headertitle}[1]{\textbf{\color{white}#1}} %titolo colonna
\definecolor{pari}{HTML}{FFDBCB}
\definecolor{dispari}{HTML}{F1F7FD}

% comandi glossario
\newcommand{\glo}{$_{G}$}
\newcommand{\glosp}{$_{G}$ }


%label custom
\makeatletter
\newcommand{\uclabel}[2]{%
	\protected@write \@auxout {}{\string \newlabel {#1}{{#2}{\thepage}{#2}{#1}{}} }%
	\hypertarget{#1}{#2}
}
\makeatother

%riportare pezzi di codice
\definecolor{codegray}{gray}{0.9}
\newcommand{\code}[1]{\colorbox{codegray}{\texttt{#1}}}



% Configurazione della pagina iniziale
\newcommand{\doctitle}{Norme di Progetto}
\newcommand{\docdate}{ } % lasciare vuoto (uno carattere di spazio) per i documenti che non sono verbali, così non viene scritta la data
\newcommand{\rev}{3.0.1}
\newcommand{\stato}{Non Approvato}
\newcommand{\uso}{Interno}
\newcommand{\approv}{--De Renzis Simone}
\newcommand{\red}{Greggio Nicolò\\& Tessari Andrea}
\newcommand{\ver}{Zuccolo Giada\\& Crivellari Alberto}
\newcommand{\dest}{Three Way Milkshake\\ Prof. Vardanega Tullio\\ Prof.
 Cardin Riccardo}
\newcommand{\describedoc}{Questo documento contiene tutte le norme di progetto\textsubscript{G}, definite inizialmente o aggiunte in seguito, viene quindi aggiornato per incrementi successivi.}



 % modifica questo file
\makeindex

\usepackage{hyperref}


\begin{document}
	\thispagestyle{empty}
\begin{titlepage}
	\begin{center}
		
		\includegraphics[scale = 0.17]{../../../../latex/images/logoTWM.png}\\[0.7cm]
		

		\noindent\rule{\textwidth}{1pt} \\[0.4cm]
		\Huge \textbf{\doctitle} \\[0.1cm]
		\ifthenelse{\equal{\docdate}{ }}{ }{ \huge \textbf{\docdate} \\[0.1cm] }
		
		\noindent\rule{\textwidth}{1pt}\\[0.7cm]
		
		\large \textbf{Three Way Milkshake - Progetto "PORTACS"} \\[0.4cm] 
                \texttt{threewaymilkshake@gmail.com} \\[0.4cm]
                
		
        
        
        \large

        \begin{tabular}{r|l}
                        \textbf{Versione} & \rev{} \\
                        \textbf{Stato} & \stato{} \\
                        \textbf{Uso} & \uso{} \\                         
                        \textbf{Approvazione} & \approv{} \\                      
                        \textbf{Redazione} & \red{} \\ 
                        \textbf{Verifica} &  \ver{} \\                         
                        \textbf{Destinatari} & \parbox[t]{5cm}{ \dest{} }
                \end{tabular} 
                \\[0.3cm]
                \large \textbf{Descrizione} \\ \describedoc{} 
               

	\end{center}
\end{titlepage}
	\pagebreak	
	
	% Registro delle modifiche
	\section*{Registro delle modifiche}

\newcommand{\changelogTable}[1]{
	
	
	\renewcommand{\arraystretch}{1.5}
	\rowcolors{2}{pari}{dispari}
	\begin{longtable}{ 
			>{\centering}p{0.07\textwidth} 
			>{}p{0.21\textwidth}
			>{\centering}p{0.17\textwidth}
			>{\centering}p{0.13\textwidth} 
			>{\centering}p{0.17\textwidth} 
			>{\centering}p{0.13\textwidth} }
		\rowcolorhead
		\headertitle{Vers.} &
		\centering \headertitle{Descrizione} &	
		\headertitle{Redazione} &
		\headertitle{Data red.} & 
		\headertitle{Verifica} &
		\headertitle{Data ver.}
		\endfirsthead	
		\endhead
		
		#1
		
	\end{longtable}
	\vspace{-2em}
	
}


\newcommand{\approvingTable}[1]{ 
	
	
	\renewcommand{\arraystretch}{1.5}
	\rowcolors{2}{pari}{dispari}
	\begin{longtable}{ 
			>{\centering}p{0.07\textwidth} 
			>{\centering}p{0.415\textwidth}
			>{\centering}p{0.13\textwidth}
			>{\centering}p{0.322\textwidth}  }
		\rowcolorhead
		\headertitle{Vers.} &
		\centering \headertitle{Descrizione} &	
		\headertitle{Data appr.} &
		\headertitle{Approvazione}
		\endfirsthead	
		\endhead
		
		#1
		
	\end{longtable}
	\vspace{-2em}
	
}
	\changelogTable{
	3.5.0 & Incremento appendice \S\ A e B & Crivellari Alberto & 2021-04-26 & -- & --
	3.4.0 & Rimossa appendice \S\ D & Crivellari Alberto & 2021-04-01 & -- & --
	\tabularnewline
	3.3.0 & Modifica sezioni \S\ 2 e \S\ 3 & Crivellari Alberto & 2021-04-01 & -- & --
	\tabularnewline
	3.2.0 & Modifica appendice \S\ B & Zuccolo Giada & 2021-03-31 & -- & --
\tabularnewline
	3.1.0 & Aggiornamento appendice \S\ C & Crivellari Alberto & 2021-03-20 & -- & --
}

\approvingTable{
    3.0.0 & Approvazione del documento & 2021-03-08 & Zuccolo Giada
}

\changelogTable{
    2.1.0 & Aggiornamento sezioni \S\ 4 e 5 & Crivellari Alberto & 2021-03-08 & Chiarello Sofia & 2021-03-08
}

\approvingTable{
    2.0.0 & Approvazione del documento & 2021-02-22 & De Renzis Simone
}

\changelogTable{
1.3.0 & Completamento stesura appendice \S\ D e \S\ B & Crivellari Alberto & 2021-02-12 & Greggio Nicolò & 2021-02-19
\tabularnewline
1.2.1 & Completamento stesura sezione \S\ 2 e \S\ 3 & Chiarello Sofia & 2021-02-11 & Tessari Andrea & 2021-02-13
\tabularnewline
1.2.0 & Stesura sezione \S\ 3 e appendice \S\ C & Chiarello Sofia & 2021-02-09 & Tessari Andrea & 2021-02-12
\tabularnewline
1.1.0 & Stesura sezione \S\ 2 & Chiarello Sofia & 2021-02-08 & Tessari Andrea & 2021-02-13
}

\approvingTable{
	1.0.0 & Approvazione del documento & 2021-01-10 & De Renzis Simone
}

\changelogTable{
0.5.1 & Aggiunta tabelle a sezione \S 4 e \S5 & Crivellari Alberto & 2021-01-08 & Greggio Nicolò & 2021-01-10
\tabularnewline
0.5.0 & Redazione sezione \S 5 & Crivellari Alberto & 2020-12-07 & Greggio Nicolò & 2021-01-10
\tabularnewline
0.4.0 & Redazione sezione \S 4 & Crivellari Alberto & 2020-12-06 & Greggio Nicolò & 2021-01-10
\tabularnewline
0.3.2 & Modifiche sezione \S 1 & Crivellari Alberto & 2020-12-30 & Greggio Nicolò & 2021-01-10
\tabularnewline
0.3.1 & Tabelle sezione \S 3 & Crivellari Alberto & 2020-12-29 & Greggio Nicolò & 2021-01-10
\tabularnewline
0.3.0 & Redazione sezione \S 3 & Tessari Andrea & 2020-12-28 & Greggio Nicolò & 2021-01-10
\tabularnewline
0.2.1 & Tabelle sezione \S 2 & Crivellari Alberto & 2020-12-20 &  Greggio Nicolò & 2021-01-09
\tabularnewline
0.2.0 & Redazione sezione \S 2 & Crivellari Alberto & 2020-12-19 & Greggio Nicolò & 2021-01-09
\tabularnewline
0.1.1 & Modifiche sezione \S 1  & Crivellari Alberto & 2020-12-18 & Greggio Nicolò & 2021-01-09
\tabularnewline
0.1.0 & Redazione sezione \S 1 & Crivellari Alberto & 2020-12-16 & Greggio Nicolò & 2021-01-09
\tabularnewline
0.0.1 & Strutturazione del documento & Tessari Andrea & 2020-12-15 & Greggio Nicolò & 2021-01-09
}


 % modifica questo file       
	\end{longtable}
	\pagebreak
		
	% indice
	\tableofcontents	
	\pagebreak
	
	% indice delle figure
	\listoffigures
	\pagebreak
	
	% indice delle tabelle
	\listoftables
	\pagebreak
		
	% contenuto del documento, ogni sezione in un file
	
	\section{Introduzione}
\subsection{Scopo del documento}
    Questo documento ha lo scopo di fissare e definire tutte le regole, convenzioni e buone pratiche utili a formare un way of working condiviso alla base da tutti i componenti del gruppo per assicurare una collaborazione efficiente ed efficace. Si discuteranno inoltre i vari strumenti che verranno adottati per facilitare lo sviluppo del progetto\textsubscript{G} e per promuovere un'organizzazione adeguata.
    Si ritiene inoltre che la definizione ed il mantenimento per incremento di un documento condiviso all'interno del gruppo di lavoro, che definisca e raccolga quanto descritto in maniera formale e centralizzata, possa favorire, in un contesto dove i membri possano variare, l'inserimento di nuovi componenti facilitandone l'ambientamento. Pur non essendo questo il contesto di lavoro attuale, è comunque una buona pratica da sperimentare e consolidare.

\subsection{Scopo del prodotto}
Il capitolato\textsubscript{G} C5 propone un progetto\textsubscript{G} in cui viene richiesto lo sviluppo di un software per il monitoraggio in tempo reale di unità che si muovono in uno spazio definito. All'interno di questo spazio, creato dall'utente per riprodurre le caratteristiche di un ambiente reale, le unità dovranno essere in grado di circolare in autonomia, o sotto il controllo dell'utente, per raggiungere dei punti di interesse posti nella mappa.  La circolazione è sottoposta a vincoli di viabilità e ad ostacoli propri della topologia dell'ambiente, deve evitare le collisioni con le altre unità e prevedere la gestione di situazioni critiche nel traffico.

\subsection{Termini, abbreviazioni ed altri documenti}
    Tutti i termini che necessitano di una spiegazione, per fornire un'adeguata comprensione, o perché possono causare ambiguità nel contesto, sono definiti nel glossario. A dispetto di tutti gli altri documenti presenti, il glossario è sotto la veste di una pagina web, più nello specifico in una wiki\textsubscript{G} di GitHub, accessibile a \href{https://github.com/Three-Way-Milkshake/docs/wiki/Glossario}{questo indirizzo}, con vocaboli suddivisi tra "Acronimi" e "Termini", disposti in ordine alfabetico al fine di facilitare la navigazione. All'interno dei documenti le voci di glossario saranno seguite da una G pedice mentre gli acronimi da una A (e.g.: voce di glossario\textsubscript{G} ; acronimo\textsubscript{A}).
    Inoltre quando si farà riferimento ad un altro documento o al documento stesso, il nome di questo sarà in maiuscoletto (e.g.: \textsc{esempio nome documento}).

\subsection{Riferimenti}
\label{ref}
    \subsubsection{Riferimenti Normativi}
        \begin{itemize}
            \item Standard ISO 12207:\\ \url{https://www.math.unipd.it/~tullio/IS-1/2009/Approfondimenti/ISO_12207-1995.pdf};
            \item Standard UML\textsubscript{A} 2.0:\\\url{https://www.omg.org/spec/UML/2.0/Superstructure/PDF};
            \item Diagrammi dei casi d'uso\textsubscript{G}:\\\url{https://www.math.unipd.it/~rcardin/swea/2021/Diagrammi\%20Use\%20Case_4x4.pdf};
            \item Diagrammi delle classi:\\\url{https://www.math.unipd.it/~rcardin/swea/2021/Diagrammi\%20delle\%20Classi_4x4.pdf};
            \item Diagrammi dei package:\\\url{https://www.math.unipd.it/~rcardin/swea/2021/Diagrammi\%20dei\%20Package_4x4.pdf};
            \item Diagrammi di attività\textsubscript{G}:\\\url{https://www.math.unipd.it/~rcardin/swea/2021/Diagrammi\%20di\%20Attivit\%c3\%a0_4x4.pdf};
            \item Diagrammi di sequenza:\\\url{https://www.math.unipd.it/~rcardin/swea/2021/Diagrammi\%20di\%20Sequenza_4x4.pdf};
            \item Standard ISO 8601:\\\url{https://www.iso.org/iso-8601-date-and-time-format.html}.
        \end{itemize}

    \subsubsection{Riferimenti Informativi}
        \begin{itemize}
        	\item \textsc{\href{https://github.com/Three-Way-Milkshake/docs/wiki/Glossario}{Glossario}}:\\per la definizione dei termini (pedice G) e degli acronimi (pedice A) evidenziati nel documento;
            \item \url{https://www.math.unipd.it/~tullio/IS-1/2020/Dispense/L03.pdf}
            \item \url{https://www.math.unipd.it/~tullio/IS-1/2020/Dispense/L06.pdf}
            \item \url{https://www.math.unipd.it/~tullio/IS-1/2020/Dispense/FC2.pdf}
            \item \url{https://www.math.unipd.it/~rcardin/swea/2021/SOLID\%20Principles\%20of\%20Object-Oriented\%20Design_4x4.pdf}
            \item \url{https://www.math.unipd.it/~tullio/IS-1/2020/Dispense/L09.pdf}
            \item \url{https://docs.github.com/en/free-pro-team@latest/actions}
            \item \url{https://www.latex-project.org/}
            \item \url{http://www.texstudio.org/}
            \item \url{https://www.xm1math.net/texmaker/}
            \item \url{https://www.atlassian.com/git/tutorials/comparing-workflows/gitflow-workflow}
            \item \url{https://github.com/about}
            \item \url{https://www.atlassian.com/software/jira}
            \item \url{https://www.atlassian.com/}
            \item \url{https://meet.google.com/}
            \item \url{https://slack.com/intl/en-it/about}
            \item \url{https://telegram.org/}
            \item \url{https://workspace.google.com/products/chat/}
        \end{itemize}

\pagebreak
 % andrea
	\pagebreak
	
	\input{res/sections/02_C1.tex} % simone
	\pagebreak

	\input{res/sections/03_C2.tex} % modifica questo file
	\pagebreak
	
	\input{res/sections/04_C3.tex} % modifica questo file
	\pagebreak
	
	\input{res/sections/05_C4.tex} % modifica questo file
	\pagebreak
	
	\input{res/sections/06_C5.tex} % modifica questo file
	\pagebreak
	
	\subsection{Capitolato C6 - RGP: Realtime Gaming Platform}
\subsubsection{Informazioni generali}
	\begin{itemize}
	\item \textbf{Nome:} RGP: Realtime Gaming Platform;
	\item \textbf{Proponente:} Zero12 s.r.l.;
	\item \textbf{Committente:}  Prof. Tullio Vardanega e Prof. Riccardo Cardin;
	\end{itemize}
\subsubsection{Descrizione del capitolato}
    Il capitolato proposto prevede la realizzazione di un videogioco a scorrimento verticale, fruibile da dispositivi mobile, con la possibilità di giocare in multiplayer real time.
    La modalità di gioco, per quanto riguarda il gameplay, sarà simile ad Aero Fighters\footnote{Esempio Aero Fighters su youtube: \url{https://youtu.be/5Phj-735p30?t=28}}, mentre la grafica potrà essere scelta liberamente dal gruppo o fornita dal team di Zero12.

\subsubsection{Finalità del progetto}
    Sono previste le modalità di gioco single e multiplayer, ma
    il cuore del progetto sono le sfide tra più giocatori: il gioco è ad eliminazione, l'ultimo rimasto vince.
    Durante la partita deve essere possibile vedere, in tempo reale, i movimenti del rivale e per garantire l'equità bisognerà prevedere la sincronizzazione di eventuali nemici e powerup.
    La modalità single player invece consiste in una serie infinita di livelli a difficoltà crescente. Il gioco termina quando l'utente ha concluso le vite oppure se non ha raccolto power-up sufficienti a mantenere il proprio oggetto attivo.
    Il progetto è finalizzato allo sviluppo di un'applicazione mobile superando dei vincoli quali:
    \begin{itemize}
    	\item ricerca delle tecnologie AWS per capire quale si può adattare meglio ad un gioco con requisiti di realtime, raccogliendo le motivazioni a supporto della scelta;
    	\item realizzazione della componente server-side;
    	\item implementazione del gioco per piattaforma mobile.
    \end{itemize}
\subsubsection{Tecnologie interessate}
    Le tecnologie interessate sono:
    \begin{itemize}
    	\item \textbf{AWS: }per la realizzazione dei servizi per la gestione dei giochi;
    	\item \textbf{NodeJs: }per lo sviluppo di codice server side;
    	\item \textbf{Swift o Kotlin: }per l'implementazione delle applicazioni mobile, rispettivamente per iOS e Android;
    	\item \textbf{Git: }per il versionamento del codice sviluppato.
    \end{itemize}
\subsubsection{Aspetti positivi}
    Il capitolato permette di entrare nel mondo dello sviluppo mobile nativo e della gestione di comunicazioni real time fra più dispositivi. Inoltre consente anche di familiarizzare con le tecnologie serverless anche tramite i servizi di Amazon AWS, per i quali l'azienda proponente prevede corsi di formazione.
\subsubsection{Criticità e fattori di rischio}
    L'ambiente AWS è molto vasto ed è critica la scelta iniziale del servizio sul quale basarsi non avendo alcuna esperienza pregressa: una scelta errata può portare ad una realizzazione non ottimale, rischiando così di non soddisfare pienamente i requisiti real time della parte multiplayer del videogioco. Inoltre lo sviluppo di applicazioni mobile native è già un compito arduo di per sé, per cui aggiungere comunicazioni real time e sincronizzazioni tra dispositivi diversi aumenta ulteriormente il grado di difficoltà e le possibilità d'errore.
\subsubsection{Conclusioni}
    I corsi di formazione offerti dall'azienda potrebbero essere interessanti, ma valutata l'elevata difficoltà e lo scarso interesse da parte del gruppo, la scelta non si pone qui. % modifica questo file
	\pagebreak
	
	\subsection{Capitolato C7 - Soluzioni di sincronizzazione Desktop}
\subsubsection{Informazioni generali}
	\begin{itemize}
	\item \textbf{Nome:} \textit{Soluzioni di sincronizzazione Desktop;}
	\item \textbf{Proponente:} \textit{Zextras;}
	\item \textbf{Committente:}  \textit{Prof. Tullio Vardanega e Prof. Riccardo Cardin;}
	\end{itemize}
\subsubsection{Descrizione del capitolato}
Il capitolato richiede di sviluppare un algoritmo di sincronizzazione Desktop solido ed efficiente in grado di garantire il salvataggio in cloud del lavoro, in modo da poter raggiungere in qualsiasi momento e da qualsiasi dispositivo il proprio lavoro. Inoltre deve essere sviluppata un'interfaccia multipiattaforma per l'uso dell'algoritmo nei più importanti sistemi operativi esistenti (MacOs, Windows, Linux) senza richiedere all'utente l'installazione di ulteriori prodotti per il funzionamento.
\subsubsection{Finalità del progetto}
L'obiettivo è quello di creare questo algoritmo di sincronizzazione che funzioni in più piattaforme in grado di garantire il salvataggio in cloud del lavoro e contemporaneamente la sincronizzazione dei cambiamenti presenti in cloud. Inoltre viene richiesto l'utilizzo dell’algoritmo sviluppato per richiedere e fornire i cambiamenti ai contenuti in sincronizzazione verso il prodotto Zextras Drive.
L'algoritmo deve avere le seguenti principali funzionalità:
\begin{itemize}
\item Configurazione ed autenticazione dell’utente;
\item Gestione di cosa sincronizzare e di cosa ignorare nelle cartelle cloud;
\item Gestione di cosa sincronizzare e di cosa ignorare nelle cartelle locali;
\item Sincronizzazione costante dei cambiamenti, siano essi locali o remoti;
\item Possibilità di modifica delle preferenze a posteriori;
\item Sistema di notifica utente dei cambiamenti.
\end{itemize}
\subsubsection{Tecnologie interessate}
L'azienda esprime la necessità di sviluppare l'algoritmo per i più importanti sistemi operativi esistenti (MacOs, Windwos, Linux) e consiglia l'utilizzo:
\begin{itemize}
\item \textbf{Qt Framework:} strumento basato sul linguaggio C++ orientato ad oggetti, da utilizzare per creare l'interfaccia poiché fortemente supportato e documentato;
\item \textbf{Python:} strumento consigliato per lo sviluppo della Business Logic, linguaggio ad alto livello adatto allo sviluppo di applicazioni distribuite.
\end{itemize}
\subsubsection{Aspetti positivi}
Il progetto si pone in un contesto molto utilizzato sia da utenti privati che da aziende. Inoltre le tecnologie consigliate fanno già parte delle conoscenze di gran parte dei membri del gruppo e ci sono molti esempi, anch'essi ben conosciuti, a cui ispirarsi per la creazione dell'algoritmo (Dropbox, Google Drive, ect).
\subsubsection{Criticità e fattori di rischio}
Il confronto con tecnologie già esistenti di questo tipo può essere svantaggioso poiché si potrebbe creare un prodotto molto più inefficiente. Inoltre le richieste sono numerose e abbastanza complesse.
\subsubsection{Conclusioni}
Si è deciso di puntare su altri capitolati poiché questo non ha suscitato grande interesse per quanto riguarda la tematica che affronta e il settore interessato. % modifica questo file
	\pagebreak	
	
	% ecc, aggiungere le sezioni successive creando i relativi file
	
\end{document}
